\documentclass{autojsarticle}

\begin{document}

\begin{center}
	\huge\bfseries autojsarticle
\end{center}

\section{概要}

これは、\LaTeX エンジンに応じて、自動的に文書クラス、文書クラスオプション、graphicxパッケージオプションを切り替えるクラスファイルです。これらの対応関係を表\ref{tbl.correspondence}に示します。

\begin{table}[h]
	\centering
	\caption{\LaTeX エンジンと文書クラス、文書クラスオプション、graphicxパッケージオプションの対応関係}
	\label{tbl.correspondence}
	\begin{tabular}{|c|c|c|c|}
		\hline 
		\LaTeX エンジン & 文書クラス & 文書クラスオプション & graphicxオプション \\ 
		\hline\hline 
		platex & jsarticle & & dvipdfmx \\ 
		\hline 
		uplatex & jsarticle & uplatex & dvipdfmx \\ 
		\hline 
		pdflatex & bxjsarticle & pdflatex,ja=standard & pdftex \\ 
		\hline 
		lualatex & ltjsarticle &  & luatex \\ 
		\hline 
		xelatex & bxjsarticle & xelatex,ja=standard & xetex \\ 
		\hline 
	\end{tabular} 
\end{table}

\section{図の貼付け}

EPS図面の貼り付け例を図\ref{fig.latex.eps}に、PDF図面の貼付け例を図\ref{fig.latex.pdf}に示します。

\begin{figure}[h]
	\centering
	\includegraphics[width=5cm]{latex.eps}
	\caption{EPS図面の貼付け例}
	\label{fig.latex.eps}
	\bigskip
	\includegraphics[width=5cm]{latex.pdf}
	\caption{PDF図面の貼付け例}
	\label{fig.latex.pdf}
\end{figure}

\section{使用法}

autojsarticle.cls を適切なディレクトリ（\slash usr\slash local\slash  texlive\slash texmf-local\slash tex\slash latex\slash local\slash 等）にコピーし、mktexlsr コマンドを実行してください。

あとは、$\backslash$documentclass\{autojsarticle\} としてやればOKです。

\verb/test/

\end{document}
